\subsection{Localization with respect to an element}

Let A be a ring and M an A-module. For every element \(f \in \mathbf{A}\) , we shall write \(\mathbf{A}_f = \mathbf{A}[f^{-1}]\) , \(\mathbf{M}_f = \mathbf{M}[f^{-1}] = \mathbf{M} \otimes_{\mathbf{A}} \mathbf{A}[f^{-1}]\) (§ 2, nos. 1 and 2); if \(\mathbf{S}_f\) is the set of \(f^n\) for \(n \geqslant 0\) , then \(\mathbf{A}_f = \mathbf{S}_f^{-1}\mathbf{A}\) , \(\mathbf{M}_f = \mathbf{S}_f^{-1}\mathbf{M}\) . If \(f\) is invertible in \(\mathbf{A}\) , \(\mathbf{A}_f\) (resp. \(\mathbf{M}_f\) ) is canonically identified with \(\mathbf{A}\) (resp. \(\mathbf{M}\) ); if \(f\) is nilpotent, then \(\mathbf{A}_f = 0\) and \(\mathbf{M}_f = 0\) . For every A-module homomorphism \(u: \mathbf{M} \to \mathbf{N}\) , we write \(u_f = u \otimes 1: \mathbf{M}_f \to \mathbf{N}_f\) .

Let \(g\) be another element of \(\mathbf{A}\) ; \(\mathbf{A}_{fg}\) (resp. \(\mathbf{M}_{fg}\) ) is canonically identified with \((\mathbf{A}_f)_{g/1}\) (resp. \((\mathbf{M}_f)_{g/1}\) ), where \(g/1\) is the image of \(g\) in \(\mathbf{A}_f\) , and \(u_{fg}\) with \((u_f)_{g/1}\) (§ 2, no. 3, Proposition 7).

PROPOSITION 1. Let \(f\) be an element of a ring \(\mathbf{A}\) and \(\phi: \mathbf{A} \to \mathbf{A}_f\) the canonical mapping. The mapping \(^a\phi: \operatorname{Spec}(\mathbf{A}_f) \to \operatorname{Spec}(\mathbf{A})\) is a homeomorphism of \(\operatorname{Spec}(\mathbf{A}_f)\) onto the open subspace \(\mathbf{X}_f\) of \(\mathbf{X} = \operatorname{Spec}(\mathbf{A})\) (§ 4, no. 3).

This a particular case of § 4, no. 3, Corollary to Proposition 13.

PROPOSITION 2. Let A be a ring, \(u: \mathbf{M} \to \mathbf{N}\) an A-module homomorphism and \(\mathfrak{p}\) a prime ideal of A.

\begin{enumerate}
\def\labelenumi{(\roman{enumi})}
\item
  Suppose that \(u_{\mathfrak{p}} \colon \mathbf{M}_{\mathfrak{p}} \to \mathbf{N}_{\mathfrak{p}}\) is surjective and that \(\mathbf{N}\) is finitely generated. Then there exists \(f \in \mathbf{A} - \mathfrak{p}\) such that \(u_{f} \colon \mathbf{M}_{f} \to \mathbf{N}_{f}\) is surjective.
\item
  Suppose that \(u_{\mathfrak{p}}\) is bijective, that \(\mathbf{M}\) is finitely generated and that \(\mathbf{N}\) is finitely presented. Then there exists \(f \in \mathbf{A} - \mathfrak{p}\) such that \(u_{f}\) is bijective.
\end{enumerate}

Let R and Q be the kernel and cokernel of u; if \(g \in A\) , the kernel and cokernel of \(u_g\) (resp. \(u_p\) ) are \(R_g\) and \(Q_g\) (resp. \(R_p\) and \(Q_p\) ) ( \(\S 2\) , no. 4, Theorem 1). Then \(Q_p = 0\) ; as N is finitely generated, so is Q and there exists \(g' \in A - p\) such that \(g'Q = 0\) ( \(\S 2\) , no. 2, Corollary 2 to Proposition 4), whence \(Q_{g'} = 0\) . Under the hypotheses of (ii), the sequence \(0 \to R_{g'} \to M_{g'} \to N_{g'} \to 0\) is exact, hence \(R_{g'}\) is finitely generated (Chapter I, \(\S 2\) , no. 8, Lemma 9). Now,

\[
\left(\mathrm{R} _ {g ^ {\prime}}\right) _ {\mathfrak{p}\mathbf{A}_{g ^ {\prime}}} = \mathrm{R} _ {\mathfrak{p}} = 0;
\]

hence there exists \(g_1 \in \mathbf{A}_{g'} - \mathfrak{p}\mathbf{A}_{g'}\) such that \(g_1\mathbf{R}_{g'} = 0\) (§ 2, no. 2, Corollary 2 to Proposition 4). Then \(g_1 = g'' / g'^h\) , where \(g'' \in \mathbf{A} - \mathfrak{p}\) ; as \(g' / 1\) is invertible in \(\mathbf{R}_{g'}\) , \((g'' / 1)\mathbf{R}_{g'} = 0\) , whence \(\mathbf{R}_{g'g''} = (\mathbf{R}_{g'})_{g'' / 1} = 0\) . If \(f = g'g''\) , \(f \in \mathbf{A} - \mathfrak{p}\) , \(\mathbf{Q}_f = 0\) and \(\mathbf{R}_f = 0\) , so that \(u_f\) is bijective.

COROLLARY. If N is finitely presented and \(N_{p}\) is a free \(A_{p}\) -module of rank p, there exists \(f \in \mathbf{A} - \mathfrak{p}\) such that \(N_{f}\) is a free \(A_{f}\) -module of rank p.

There exist by hypothesis \(p\) elements \(x_{i} \in \mathbf{N}\) ( \(1 \leqslant i \leqslant p\) ) such that the \(x_{i}/1\) form a basis of the free \(\mathbf{A}_{\mathfrak{p}}\) -module \(\mathbf{N}_{\mathfrak{p}}\) . Consider the homomorphism \(u: \mathbf{A}^p \to \mathbf{N}\) such that \(u(e_{i}) = x_{i}\) for \(1 \leqslant i \leqslant p\) , \((e_{i})_{1 \leqslant i \leqslant p}\) being the canonical basis of \(\mathbf{A}^p\) . As \(u_{\mathfrak{p}}\) is bijective by hypothesis, there exists \(f \in \mathbf{A} - \mathfrak{p}\) such that \(u_{f}\) is bijective, by virtue of Proposition 2.

PROPOSITION 3. Let \((f_{i})_{i\in I}\) be a finite family of elements of a ring A, generating the ideal A of A. The ring \(\mathbf{B} = \prod_{i\in I}\mathbf{A}_{f_i}\) is then a faithfully flat A-module.

By §2, no. 4, Theorem 1, each of the \(\mathbf{A}_{f_i}\) is a flat A-module, hence so is B (Chapter I, §2, no. 3, Proposition 2). On the other hand, if \(\mathfrak{p}\) is a prime ideal of A, there exists an index \(i\) such that \(f_i \notin \mathfrak{p}\) and \(\mathfrak{p}_{f_i} = \mathfrak{p} \mathbf{A}_{f_i}\) is therefore a prime ideal of \(\mathbf{A}_{f_i}\) . Then \(\mathfrak{p} \mathbf{B} \subset \mathfrak{p} \mathbf{A}_{f_i} \times \prod_{j \neq i} \mathbf{A}_{f_j} \neq \mathbf{B}\) since \(\mathfrak{p} \mathbf{A}_{f_i} \neq \mathbf{A}_{f_i}\) ; this suffices to imply that B is a faithfully flat A-module (Chapter I, §3, no. 1, Proposition 1).

COROLLARY. Under the hypotheses of Proposition 3, for an A-module M to be finitely generated (resp. finitely presented), it is necessary and sufficient that, for every index i, the \(A_{f_{i}}\) -module \(M_{f_{i}}\) be finitely generated (resp. finitely presented).

The condition is obviously necessary (§ 2, no. 4). Conversely, if all the \(M_{f_{i}}\) are finitely generated (resp. finitely presented), \(M' = \prod_{i \in I} M_{f_{i}}\) is a finitely generated (resp. finitely presented) B-module, for we can obviously suppose that for each i there is an exact sequence \(A_{f_i}^{m} \to A_{f_i}^{n} \to M_{f_i} \to 0\) , where m and n are independent of i). Now, \(M' = M \otimes_{A} B\) . The corollary then follows from Proposition 3 and Chapter I, § 3, no. 6, Proposition 11.

Note that the condition on the \(f_{i}\) means that the open sets \(\mathbf{X}_{f_{i}}\) form a covering of \(\operatorname{Spec}(\mathbf{A})\) ( \(\S 4\) , no. 3, Corollary 3 to Proposition 11).

